\section*{Chapter \ref{ch:b2:exchange-economies} --- Exchange Economies}

\begin{solution}{exo:b2:exchange-economies:1}
% code: exercises/solutions.py ex_feasible
(i) Totals $(10, 10)$: \omterm{def:b2:exchange-economies:feasible}{feasible} and non-wasteful. (ii) Totals $(11, 9)$: not \omterm{def:b2:exchange-economies:feasible}{feasible}, since 11 units of
good 1 are consumed and only 10 exist. (iii) Totals $(9, 9)$: \omterm{def:b2:exchange-economies:feasible}{feasible} but wasteful; one unit of each
good is thrown away.
\end{solution}

\begin{solution}{exo:b2:exchange-economies:2}
% code: exercises/solutions.py ex_mrs
$\mathit{MRS}_A = x_{A2}/x_{A1} = 6/4 = 1.5$ and $\mathit{MRS}_B = x_{B2}/(2x_{B1}) = 4/12 = 0.333$. They
differ, so the \omterm{def:b2:exchange-economies:allocation}{allocation} is not efficient. $A$ values good 1 at 1.5 units of good 2, $B$ at a third of a
unit: good 1 should move from $B$ to $A$, against good 2, at any rate between the two.
\end{solution}

\begin{solution}{exo:b2:exchange-economies:3}
% figdata: smith_wide (T6A 49.2, T4A 107.1, T9A 90.3)
Test 6A started furthest away, at 49.2 per cent of the prediction. In the last period, test 4A was still
7.1 per cent above (107.1) and test 9A 9.7 per cent below (90.3); the others ended within 2 per cent.
Tests 1 and 6A cross the prediction from below, in their third and fourth periods.
\end{solution}

\begin{solution}{exo:b2:exchange-economies:4}
Many markets with the same design in which prices stay away from the prediction period after period, or
in which the dispersion $\alpha$ does not fall, or in which prices settle at a level that depends on the
first trades rather than on the reservation prices. One period proves nothing: every test of the chart
began more than 5 per cent away from the prediction, and most then converged.
\end{solution}

\begin{solution}{exo:b2:exchange-economies:5}
% code: exercises/solutions.py ex_cd_eq
$a = \tfrac12$, $b = \tfrac14$, so $p^* = (\tfrac12 \times 2 + \tfrac14 \times 8)/(\tfrac12 \times 6 +
\tfrac34 \times 4) = 3/6 = 0.5$. $A$'s income is $0.5 \times 6 + 2 = 5$: $x_A = (0.5 \times 5/0.5,\ 0.5
\times 5) = (5, 2.5)$. $B$'s income is $0.5 \times 4 + 8 = 10$: $x_B = (0.25 \times 10/0.5,\ 0.75 \times
10) = (5, 7.5)$. Totals: $5 + 5 = 10$ and $2.5 + 7.5 = 10$.
\end{solution}

\begin{solution}{exo:b2:exchange-economies:6}
% code: exercises/solutions.py ex_contract
By \eqref{eq:b2:exchange-economies:contract}, $x_{A2} = \tfrac14 \cdot \tfrac12 \cdot 10\,x_{A1}/(\tfrac12
\cdot \tfrac34 (10 - x_{A1}) + \tfrac14 \cdot \tfrac12 x_{A1}) = 5x_{A1}/(15 - x_{A1})$. At $x_{A1} = 2$,
$5$, $8$: $0.769$, $2.5$, $5.714$. The equilibrium $(5, 2.5)$ is on it. The curve lies below the
diagonal because $A$ spends a larger share on good 1 than $B$ does, so efficient \omterm{def:b2:exchange-economies:allocation}{allocations} give $A$
relatively more of good 1.
\end{solution}

\begin{solution}{exo:b2:exchange-economies:7}
% code: exercises/solutions.py ex_transfer
$A$ spends half of any extra income on good 1, $B$ only a quarter: the transfer raises world demand for
good 1, so $p$ rises. New \omterm{def:b2:exchange-economies:endowment}{endowments} $(6, 4)$ and $(4, 6)$: $p^* = (\tfrac12 \times 4 + \tfrac14 \times
6)/6 = 3.5/6 = 0.5833$, up by $0.0833$. In the formula for $p^*$ the denominator does not involve good 2,
and the numerator changes by $a\Delta - b\Delta$, so $\Delta p^* = (a - b)\Delta/((1 - a)\omega_{A1} + (1
- b)\omega_{B1}) = 0.25 \times 2/6 = 0.0833$.
\end{solution}

\begin{solution}{exo:b2:exchange-economies:8}
% code: exercises/solutions.py ex_walras
At $p = 1$, $A$'s income is 8 and $x_A = (4, 4)$; $B$'s is 12 and $x_B = (3, 9)$. So $z_1 = 7 - 10 = -3$
and $z_2 = 13 - 10 = 3$, and $1 \times (-3) + 1 \times 3 = 0$. Good 1 is in excess supply: its relative
price $p$ must fall, towards the equilibrium $0.5$.
\end{solution}

\begin{solution}{exo:b2:exchange-economies:9}
Induce the same values and costs, with two equilibrium prices, in many sessions with fresh participants,
and vary only where trading starts: for instance, the first period's trades are made at a price imposed
near one equilibrium in some sessions and near the other in the rest, after which trading is free. If
sessions started near the low equilibrium settle there and those started near the high one settle there,
history picks the equilibrium. If all settle at one of them, something outside the model (stability, a
focal point) selects it; the model, which predicts only the set, is not refuted by either result.
\end{solution}

\begin{solution}{exo:b2:exchange-economies:10}
Local non-satiation means each consumer spends her whole income: $\sum_\ell p_\ell x_{i\ell}(p) =
\sum_\ell p_\ell\omega_{i\ell}$ (a bundle that left income unspent could be improved by a nearby bundle
still within the budget). Summing over $i$: $\sum_\ell p_\ell z_\ell(p) = 0$. If $z_\ell(p) = 0$ for
$\ell = 1, \dots, n - 1$, then $p_n z_n(p) = 0$, and $p_n > 0$ gives $z_n(p) = 0$. Example: a third good
that no consumer values, held in positive amounts. Preferences are still locally non-satiated through
the other goods. At any equilibrium its price is zero (at a positive price everyone would sell it and no
one buy it), and it is in excess supply: $z_3 = -\bar\omega_3 < 0$, while $p_3z_3 = 0$ keeps \omterm{def:b2:exchange-economies:walras-law}{Walras' law}.
\end{solution}

\begin{solution}{exo:b2:exchange-economies:11}
% code: exercises/solutions.py three
(a) $A$ maximises $x_{A1} - \tfrac18 x_{A2}^{-8}$ subject to $px_{A1} + x_{A2} = 2p + r$. The ratio of
marginal utilities $x_{A2}^{-9}/1$ equals the price ratio $1/p$, so $x_{A2} = p^{1/9}$ and $x_{A1} = 2 +
(r - p^{1/9})/p$. $B$ maximises $-\tfrac18x_{B1}^{-8} + x_{B2}$: $x_{B1}^{-9} = p$, so $x_{B1} = p^{-1/9}$.
Then $z_1(p) = x_{A1} + x_{B1} - (2 + r) = (r - p^{1/9})/p + p^{-1/9} - r$.
(b) $z_1(1) = (r - 1) + 1 - r = 0$. $z_1(2) = (r - 2^{1/9})/2 + 2^{-1/9} - r = -\tfrac12 r - \tfrac12
2^{1/9} + 2^{-1/9}$, and with $r = 2^{8/9} - 2^{1/9}$, $\tfrac12 r = 2^{-1/9} - \tfrac12 2^{1/9}$: the sum
is zero. $z_1(\tfrac12) = 2(r - 2^{-1/9}) + 2^{1/9} - r = r - 2^{8/9} + 2^{1/9} = 0$.
(c) Numerically, $z_1'(\tfrac12) = -0.0348$, $z_1'(1) = 0.0061$, $z_1'(2) = -0.0043$. Under $\dot p =
kz_1(p)$ an equilibrium is locally stable when $z_1' < 0$: $p = \tfrac12$ and $p = 2$ are stable, $p = 1$
is not, as the listing showed.
\end{solution}

\begin{solution}{exo:b2:exchange-economies:12}
Feasibility in the two-consumer economy gives $x_A + x_B = \omega_A + \omega_B$, so $x_B - \omega_B =
\omega_A - x_A$. Let the coalition $\{A_1, A_2, B_1\}$ give $B_1$ the bundle $x_B$, and each $A$ consumer
$y = \tfrac12(\omega_A + x_A)$. This is \omterm{def:b2:exchange-economies:feasible}{feasible} for the coalition: $2y + x_B = \omega_A + x_A + x_B =
2\omega_A + \omega_B$, its own \omterm{def:b2:exchange-economies:endowment}{endowments}. $B_1$ is exactly as well off as under the \omterm{def:b2:exchange-economies:allocation}{allocation}. Each $A$
consumer is strictly better off: $x_A$ and $\omega_A$ are different bundles on the same indifference
curve, and strict quasi-concavity makes their midpoint strictly preferred to both. So the coalition
blocks. The \omterm{def:b2:exchange-economies:allocation}{allocation} is in the \omterm{def:b2:exchange-economies:core}{core} of the two-consumer economy (it is efficient and both weakly
prefer it to their \omterm{def:b2:exchange-economies:endowment}{endowments}) but not in the \omterm{def:b2:exchange-economies:core}{core} of the replica: replication removed an end of the
\omterm{def:b2:exchange-economies:core}{core}.
\end{solution}

\begin{solution}{pb:b2:exchange-economies:1}
% code: exercises/solutions.py problem
\textbf{Part I.}
1. The box is 80 tonnes of fish wide and 100 of grain high; the \omterm{def:b2:exchange-economies:endowment}{endowment} is $(60, 10)$ from Brenland's
corner, $(20, 90)$ from Kestria's.
2. Brenland: $\mathit{MRS} = (\tfrac12/\tfrac12)(10/60) = 0.1667$ tonnes of grain per tonne of fish.
Kestria: $(\tfrac14/\tfrac34)(90/20) = 1.5$.
3. No. Any small trade of fish from Brenland to Kestria at between $0.1667$ and $1.5$ tonnes of grain per
tonne of fish makes both better off.
4. Brenland moves to $(50, 15)$: $\sqrt{50 \times 15} = 27.39 > \sqrt{60 \times 10} = 24.49$. Kestria
moves to $(30, 85)$: $30^{1/4}85^{3/4} = 65.52 > 20^{1/4}90^{3/4} = 61.79$. Both gain.
5. Brenland exports fish because, holding much fish and little grain, it values a tonne of fish far less
in grain than Kestria does.
\textbf{Part II.}
6. $\mathit{MRS}_B = \mathit{MRS}_K$: $g_B/f_B = \tfrac13(100 - g_B)/(80 - f_B)$.
7. By \eqref{eq:b2:exchange-economies:contract} with $a = \tfrac12$, $b = \tfrac14$: $g_B = 12.5f_B/(30 -
0.25f_B) = 50f_B/(120 - f_B)$.
8. $f_B = 0$ gives $g_B = 0$ and $f_B = 80$ gives $4000/40 = 100$. At $f_B = 40$, $g_B = 25$, below the
diagonal's $50$: Brenland spends a larger share on fish, so efficient \omterm{def:b2:exchange-economies:allocation}{allocations} give it relatively
more fish.
9. From $f_B = 32.42$ ($g_B = 18.51$, where Brenland is as well off as at its \omterm{def:b2:exchange-economies:endowment}{endowment}) to $f_B = 42.25$
($g_B = 27.17$, where Kestria is).
10. Yes: any change takes something from Brenland. Efficiency rules out waste; it does not rule out an
\omterm{def:b2:exchange-economies:allocation}{allocation} that leaves one country with nothing.
\textbf{Part III.}
11. Brenland: $f_B = \tfrac12(60p + 10)/p$, $g_B = \tfrac12(60p + 10)$. Kestria: $f_K = \tfrac14(20p +
90)/p$, $g_K = \tfrac34(20p + 90)$.
12. $z_f(p) = (30p + 5 + 5p + 22.5)/p - 80 = 27.5/p - 45$.
13. $p^* = 27.5/45 = 0.6111$ tonnes of grain per tonne of fish.
14. Brenland's income is $0.6111 \times 60 + 10 = 46.67$: $f_B = 38.18$, $g_B = 23.33$. Kestria's is
$0.6111 \times 20 + 90 = 102.22$: $f_K = 41.82$, $g_K = 76.67$. Brenland exports $21.82$ tonnes of fish
and imports $13.33$ of grain ($21.82 \times 0.6111 = 13.33$).
15. $23.33 + 76.67 = 100$: the grain market clears, as \omterm{def:b2:exchange-economies:walras-law}{Walras' law} requires once the fish market does.
16. $50 \times 38.18/(120 - 38.18) = 23.33$, and $38.18$ lies between $32.42$ and $42.25$.
17. Brenland: $\sqrt{38.18 \times 23.33} = 29.85$ against $24.49$. Kestria: $41.82^{1/4}76.67^{3/4} =
65.89$ against $61.79$. Both gain.
18. Brenland: $60(10 + E) = 38.18 \times 23.33 = 890.7$, so $10 + E = 14.85$ and $E = 4.85$ tonnes of
grain. Kestria: $20^{1/4}(90 + E)^{3/4} = 41.82^{1/4}76.67^{3/4}$ gives $90 + E = 76.67 \times
(41.82/20)^{1/3} = 76.67 \times 1.2787 = 98.04$, so $E = 8.04$ tonnes.
\textbf{Part IV.}
19. The price of fish rises: Brenland spends half of any income on fish, Kestria a quarter, so moving
income from Kestria to Brenland raises the world demand for fish.
20. \omterm{def:b2:exchange-economies:endowment}{Endowments} $(60, 20)$ and $(20, 80)$: $p^* = (\tfrac12 \times 20 + \tfrac14 \times 80)/45 = 30/45 =
0.6667$. Brenland's income is $0.6667 \times 60 + 20 = 60$: $(f_B, g_B) = (45, 30)$. Kestria's is
$0.6667 \times 20 + 80 = 93.33$: $(35, 70)$.
21. Brenland gains: $\sqrt{45 \times 30} = 36.74 > 29.85$. Kestria loses: $35^{1/4}70^{3/4} = 58.86 <
65.89$. It still gains from trade given its new \omterm{def:b2:exchange-economies:endowment}{endowment}, which is worth $20^{1/4}80^{3/4} = 56.57 <
58.86$.
22. $50 \times 45/(120 - 45) = 30$: on the same curve.
23. Moving \omterm{def:b2:exchange-economies:endowment}{endowments} moves the equilibrium along the \omterm{def:b2:exchange-economies:contract-curve}{contract curve}; the next chapter shows that, with
convex preferences, any efficient \omterm{def:b2:exchange-economies:allocation}{allocation} can be reached this way.
24. Fish trades at $0.6111$ tonnes of grain per tonne; trade is worth $4.85$ tonnes of grain to Brenland
and $8.04$ to Kestria.
\end{solution}
